\Needspace{13cm}
\subsection{4.3 核心对象状态机}
\noindent 图上边编号对照随后转移表。每条转移按“触发；[守卫]；执行者；副作用”解释；守卫失败保持原状态并记拒绝原因。订单、单件、备货、鉴定、维修、押金独立建模，防止一个状态字段混合多种责任。订单图中的 O8 虚线为按服务器时间推导的生效读取状态，并非持久写入。
% 边上的 O/U/P 编号对照 analysis-states.tex 的触发/守卫/岗位/副作用表。
\newcommand{\BSRStateOrder}{\begin{tikzpicture}[x=1cm,y=1cm]
  \node[bsrState] (draft) at (0,0) {草稿};
  \node[bsrState] (hold) at (3.2,0) {暂留};
  \node[bsrState] (confirmed) at (6.4,0) {已确认};
  \node[bsrState] (prepare) at (9.6,0) {备货中};
  \node[bsrState] (renting) at (12.8,0) {租赁中};
  \node[bsrState] (cancel) at (3.2,-3) {已取消};
  \node[bsrState] (expired) at (0,-3) {暂留已过期};
  \node[bsrState] (settle) at (9.6,-3) {待结算};
  \node[bsrState] (done) at (12.8,-3) {完成};
  \draw[bsrMessage] (draft) -- node[above,font=\tiny]{O1} (hold);
  \draw[bsrMessage] (hold) -- node[above,font=\tiny]{O2} (confirmed);
  \draw[bsrMessage] (confirmed) -- node[above,font=\tiny]{O3} (prepare);
  \draw[bsrMessage] (prepare) -- node[above,font=\tiny]{O4} (renting);
  \draw[bsrMessage] (renting) -- node[right,font=\tiny]{O5} (settle);
  \draw[bsrMessage] (settle) -- node[above,font=\tiny]{O6} (done);
  \draw[bsrMessage] (hold) -- node[right,font=\tiny]{O7} (cancel);
  \draw[bsrMessage,dashed] (hold) to[out=220,in=50] node[above,font=\tiny]{O8} (expired);
\end{tikzpicture}}

\newcommand{\BSRStateUnit}{\begin{tikzpicture}[x=1cm,y=1cm]
  \node[bsrState] (available) at (0,0) {可租};
  \node[bsrState] (reserved) at (3.3,0) {准备占用};
  \node[bsrState] (out) at (6.6,0) {借出};
  \node[bsrState] (dirty) at (9.9,0) {待清洁};
  \node[bsrState] (isolate) at (6.6,-3) {疑损隔离};
  \node[bsrState] (repair) at (9.9,-3) {维修};
  \node[bsrState] (accept) at (13.2,-3) {待安全验收};
  \node[bsrState] (retired) at (13.2,-5.4) {报废};
  \draw[bsrMessage] (available) -- node[above,font=\tiny]{U1} (reserved);
  \draw[bsrMessage] (reserved) -- node[above,font=\tiny]{U2} (out);
  \draw[bsrMessage] (out) -- node[above,font=\tiny]{U3} (dirty);
  \draw[bsrMessage] (dirty) to[out=90,in=90,looseness=1] node[above,font=\tiny]{U4} (available);
  \draw[bsrMessage] (out) -- node[left,font=\tiny]{U5} (isolate);
  \draw[bsrMessage] (isolate) -- node[above,font=\tiny]{U6} (repair);
  \draw[bsrMessage] (repair) -- node[above,font=\tiny]{U7} (accept);
  \draw[bsrMessage] (accept) to[out=70,in=25,looseness=1] node[above,font=\tiny]{U8} (available);
  \draw[bsrMessage] (accept) -- node[right,font=\tiny]{U9} (retired);
\end{tikzpicture}}

\newcommand{\BSRStatePreparation}{\begin{tikzpicture}[x=1cm,y=1cm]
  \node[bsrState] (new) at (0,0) {待分派};
  \node[bsrState] (assigned) at (3.4,0) {已指派};
  \node[bsrState] (picking) at (6.8,0) {逐件准备};
  \node[bsrState] (ready) at (10.2,0) {已备妥};
  \node[bsrState] (handed) at (13.6,0) {已交付};
  \node[bsrState] (blocked) at (6.8,-3) {阻塞/替换};
  \draw[bsrMessage] (new) -- node[above,font=\tiny]{P1} (assigned);
  \draw[bsrMessage] (assigned) -- node[above,font=\tiny]{P2} (picking);
  \draw[bsrMessage] (picking) -- node[above,font=\tiny]{P3} (ready);
  \draw[bsrMessage] (ready) -- node[above,font=\tiny]{P4} (handed);
  \draw[bsrMessage] (picking) -- node[right,font=\tiny]{P5} (blocked);
  \draw[bsrMessage] (blocked) to[out=180,in=220,looseness=1] node[left,font=\tiny]{P6} (picking);
\end{tikzpicture}}

% 边上的 D/M/F 编号对照 analysis-states.tex。
\newcommand{\BSRStateDamage}{\begin{tikzpicture}[x=1cm,y=1cm]
  \node[bsrState] (suspect) at (0,0) {疑似受损};
  \node[bsrState] (assigned) at (3.4,0) {技师已指派};
  \node[bsrState] (initial) at (6.8,0) {24h 初评};
  \node[bsrState] (review) at (10.2,0) {经理复核};
  \node[bsrState] (approved) at (13.6,0) {核定};
  \node[bsrState] (supplement) at (6.8,-3) {退回补证};
  \node[bsrState] (dispute) at (13.6,-3) {争议};
  \draw[bsrMessage] (suspect) -- node[above,font=\tiny]{D1} (assigned);
  \draw[bsrMessage] (assigned) -- node[above,font=\tiny]{D2} (initial);
  \draw[bsrMessage] (initial) -- node[above,font=\tiny]{D3} (review);
  \draw[bsrMessage] (review) -- node[above,font=\tiny]{D4} (approved);
  \draw[bsrMessage] (review) -- node[sloped,above,font=\tiny]{D5} (supplement);
  \draw[bsrMessage] (supplement) -- node[right,font=\tiny]{D6} (initial);
  \draw[bsrMessage] (review) -- node[right,font=\tiny]{D7} (dispute);
\end{tikzpicture}}

\newcommand{\BSRStateMaintenance}{\begin{tikzpicture}[x=1cm,y=1cm]
  \node[bsrState] (queued) at (0,0) {待接单};
  \node[bsrState] (diagnose) at (3.4,0) {诊断};
  \node[bsrState] (repair) at (6.8,0) {维修};
  \node[bsrState] (review) at (10.2,0) {待安全验收};
  \node[bsrState] (safe) at (13.6,0) {可租};
  \node[bsrState] (parts) at (6.8,-3) {待配件};
  \node[bsrState] (reject) at (10.2,-3) {退回重修};
  \node[bsrState] (scrap) at (13.6,-3) {报废};
  \draw[bsrMessage] (queued) -- node[above,font=\tiny]{M1} (diagnose);
  \draw[bsrMessage] (diagnose) -- node[above,font=\tiny]{M2} (repair);
  \draw[bsrMessage] (repair) -- node[above,font=\tiny]{M3} (review);
  \draw[bsrMessage] (repair) -- node[right,font=\tiny]{M4} (parts);
  \draw[bsrMessage] (parts) to[out=180,in=220,looseness=1] node[left,font=\tiny]{M5} (repair);
  \draw[bsrMessage] (review) -- node[above,font=\tiny]{M6} (safe);
  \draw[bsrMessage] (review) -- node[right,font=\tiny]{M7} (reject);
  \draw[bsrMessage] (reject) -- node[above,font=\tiny]{M8} (repair);
  \draw[bsrMessage] (review) -- node[sloped,above,font=\tiny]{M9} (scrap);
\end{tikzpicture}}

\newcommand{\BSRStateDeposit}{\begin{tikzpicture}[x=1cm,y=1cm]
  \node[bsrState] (hold) at (0,0) {押金暂持};
  \node[bsrState] (calc) at (3.4,0) {费用待核};
  \node[bsrState] (approved) at (6.8,0) {结算已核};
  \node[bsrState] (refunding) at (10.2,0) {退款发起};
  \node[bsrState] (done) at (13.6,0) {已到账/完成};
  \node[bsrState] (dispute) at (3.4,-3) {争议冻结};
  \node[bsrState] (supplement) at (6.8,-3) {待联系补差};
  \node[bsrState] (failed) at (10.2,-3) {退款失败/待核实};
  \draw[bsrMessage] (hold) -- node[above,font=\tiny]{F1} (calc);
  \draw[bsrMessage] (calc) -- node[above,font=\tiny]{F2} (approved);
  \draw[bsrMessage] (approved) -- node[above,font=\tiny]{F3} (refunding);
  \draw[bsrMessage] (refunding) -- node[above,font=\tiny]{F4} (done);
  \draw[bsrMessage] (calc) -- node[right,font=\tiny]{F5} (dispute);
  \draw[bsrMessage] (approved) -- node[right,font=\tiny]{F6} (supplement);
  \draw[bsrMessage] (refunding) -- node[right,font=\tiny]{F7} (failed);
  \draw[bsrMessage] (failed) to[out=55,in=-55,looseness=1.2] node[right,font=\tiny]{F8} (refunding);
\end{tikzpicture}}

\BSRNewFigure{\BSRStateOrder}{ST-ORDER 订单状态：暂留、履约和最终结清}{fig:st-order}
\BSRNewFigure{\BSRStateUnit}{ST-UNIT 物理单件状态：隔离、维修与经理验收}{fig:st-unit}
\BSRNewFigure{\BSRStatePreparation}{ST-PREP 备货/交付状态：具名处理与阻塞替换}{fig:st-prep}
\BSRNewFigure{\BSRStateDamage}{ST-DAMAGE 受损鉴定：24 小时初评与争议}{fig:st-damage}
\BSRNewFigure{\BSRStateMaintenance}{ST-MAINT 维修工单：技师提交不等于可租}{fig:st-maint}
\BSRNewFigure{\BSRStateDeposit}{ST-DEPOSIT 押金与退款：争议、补差和失败回执}{fig:st-deposit}

\subsubsection{4.3.1 订单与单件转移}
\begin{BSRTable}
\rowcolor{TableHead}\color{white}\bfseries 边 & \color{white}\bfseries 触发；[守卫]；执行者；副作用 \\
O1 & 提交预约；[规则同意、报价新鲜、整段有货]；顾客/系统；原子建暂留与截止。\\
O2 & 确认付款；[当前合同版签署、金额回执有效；演示另标]；支付控制；冻结订单金额。\\
O3 & 付款成功；[仅一笔有效交易]；系统；创建一份备货任务并指派门店。\\
O4 & 完成交付；[当前版签署、核身份码、单件安全检查通过]；专员；逐件写借出和责任人。\\
O5 & 最后一件归还；[所有实还事实齐]；专员/系统；进入待鉴定与结算。\\
O6 & 结清；[费用核定、无争议、补差/退款回执闭环]；经理/系统；订单完成。\\
O7 & 顾客取消；[仍在有效暂留，或已确认且符合取消规则]；顾客；持久记取消、费用及审计。\\
O8 & 暂留到期；[原始 HOLD 且服务器时间达到截止]；系统读取；返回 EXPIRED、释放库存和提醒占用，旧行及报价不改写。\\
U1 & 绑定备货；[整段未重叠、所在店正确、可租]；专员/系统；单件占用。\\
U2 & 逐件交付；[安全和合同/取货码守卫通过]；专员；记借出时刻。\\
U3 & 归还事实；[单件未归还、检查为需清洁]；接收专员；冻结实还并改当前店。\\
U4 & 清洁通过；[检查合格且无开放工单]；专员；恢复可订量。\\
U5 & 发现疑损；[有实物与订单对应]；接收专员；隔离并派鉴定。\\
U6 & 开维修工单；[已隔离且无重复开放工单]；经理/专员；冻结技师费率。\\
U7 & 技师提交；[诊断、工时、安全项完整]；技师；转待经理验收，仍不可租。\\
U8 & 安全验收通过；[经理复核、无其他开放工单、已清洁]；经理；恢复可租。\\
U9 & 判报废；[经理有依据确认]；经理；永久退出可订量并留估值记录。\\
\end{BSRTable}

\subsubsection{4.3.2 备货与受损转移}
\begin{BSRTable}
\rowcolor{TableHead}\color{white}\bfseries 边 & \color{white}\bfseries 触发；[守卫]；执行者；副作用 \\
P1 & 指派任务；[已确认订单且处理人有门店权限]；经理/系统；顾客见具名责任人。\\
P2 & 接单拣货；[专员在岗且库存行未取消]；专员；进入逐件绑定。\\
P3 & 备货完成；[每件规格/尺码/安全项通过]；专员；顾客见已备妥。\\
P4 & 核码交付；[身份、码、当前合同版有效]；专员；逐件写借出。\\
P5 & 发现不合格/占用；[绑定或检查失败]；专员；隔离问题件、记录阻塞。\\
P6 & 替换再拣货；[同规格替代件可租]；专员/系统；释放旧绑定并重新检查。\\
D1 & 发现疑损；[逐件归还事实已保存]；系统/经理；隔离并指派一名技师。\\
D2 & 提交初评；[部位照片、事实和级别齐全]；技师；记录发现后 24h 是否及时并告知顾客。\\
D3 & 请求复核；[初评版本已锁定]；技师；经理收到待办，不直接扣费。\\
D4 & 核定；[证据充分、正常磨损零费]；经理；冻结损坏费依据。\\
D5 & 退回补证；[照片/估值不足]；经理；保留旧版本并继续隔离。\\
D6 & 补证重交；[新增证据完整]；技师；生成新初评版本，不覆写旧照片。\\
D7 & 顾客提出异议；[案件未最终结清]；顾客/经理；标争议并暂停押金结清。\\
\end{BSRTable}

\subsubsection{4.3.3 维修与押金转移}
\begin{BSRTable}
\rowcolor{TableHead}\color{white}\bfseries 边 & \color{white}\bfseries 触发；[守卫]；执行者；副作用 \\
M1 & 接单；[本人获指派且单件已隔离]；技师；记录责任与时间。\\
M2 & 诊断后维修；[故障确认]；技师；登记工时与材料。\\
M3 & 提交安全验收；[工时为 0.25 小时倍数、四项检查齐]；技师；仍保持维修隔离。\\
M4 & 等待配件；[有缺件说明]；技师；暂停工作但保留隔离。\\
M5 & 配件到位；[获授权记录]；技师；继续维修。\\
M6 & 经理验收通过；[实物安全/清洁/无其他开放单]；经理；恢复可租。\\
M7 & 经理退回；[安全项不合格或证据不足]；经理；保留诊断并回到重修。\\
M8 & 重新处理；[技师接收退回意见]；技师；追加工作事件。\\
M9 & 报废；[经理确认不可修复]；经理；单件永久退出可订。\\
F1 & 最后一件实还；[逐件借还完整]；系统；进入费用待核。\\
F2 & 核定费用；[损坏初评复核完成、无争议]；经理；冻结超时/损坏明细。\\
F3 & 发起退款；[补差为零且真实/演示路径已标]；系统；写幂等资金请求。\\
F4 & 收到已到账回执；[可验证第三方回执；演示不得冒充]；支付服务/系统；标完成。\\
F5 & 顾客争议；[未最终结清]；顾客/经理；冻结退款和补差终结。\\
F6 & 费用超押金；[补差大于零]；系统/专员；派员工联系，不自动追扣。\\
F7 & 退款失败/回执不明；[外部响应失败或超时]；系统；标失败/待核实并提醒经理。\\
F8 & 核对后重试；[经理核实原请求无成功入账]；经理/支付控制；沿原幂等引用重试。\\
\end{BSRTable}

\noindent \textbf{跨对象不变量：}“提醒已读”仅清除本人新增提示，不关闭业务案件；“技师已提交”不改变设备可租；“部分归还”不使押金自动结清；目标真实支付与本地模拟事件不得混同。
